% ============================================================
% 三、数学背景与直觉 (Math Background & Intuition)
% ============================================================
\section{数学背景与直觉}

% 提示：本节的核心目标是"用一句话、一张图说明它究竟在做什么"，
% 而不是堆砌公式。多用自己的话和具体的小例子。

\subsection{核心概念概述}

列出本文涉及的核心数学概念，每个概念用一段话说明它"是什么"以及"为什么在此处有用"。

\begin{itemize}
    \item \textbf{【概念 1，例如：梯度下降】}：【用自己的话解释梯度下降的几何意义：每次沿局部最陡下降方向移动，就像蒙着眼睛下山时每次朝脚底最陡的方向迈一步。关联课程第 X 章。】
    \item \textbf{【概念 2，例如：对偶问题】}：【用自己的话解释对偶的物理/几何直觉……关联课程第 Y 章。】
    \item \textbf{【概念 3，…】}：【……】
\end{itemize}

\subsection{直观解释}

\subsubsection{一个小例子}

用一个极度简化的小规模例子（2--3 维，手算级别）展示核心方法到底在做什么。

\begin{example}
    考虑一个二维问题：
    \begin{equation}
        \min_{x_1, x_2} \; f(x_1, x_2) = (x_1 - 1)^2 + 10(x_2 - 2)^2
    \end{equation}
    这是一个典型的病态二次函数……【解释梯度下降在此处的行为：zigzag 现象及其几何原因】
\end{example}

\subsubsection{几何/物理图像}

% === 建议在此处插入一张手绘示意图或矢量图 ===

\begin{figure}[htbp]
    \centering
    % \includegraphics[width=0.6\textwidth]{figures/intuition.pdf}
    \caption{【图的标题：例如，梯度下降在病态二次函数上的 zigzag 轨迹（等高线图）】}
    \label{fig:intuition}
\end{figure}

图~\ref{fig:intuition} 展示了……【解释图中各元素的含义】

\subsection{与课程内容的关联}

\begin{itemize}
    \item 本文的核心数学工具对应课程讲义第【X】章（【具体主题】）；
    \item 此外，推导中用到第【Y】章的【具体知识点】；
    \item 【如果有】课程中未涵盖但本文用到的额外数学知识，简要说明并给出参考文献。
\end{itemize}
